\chapter{Lewis Structures, Resonance and VSEPR}\label{ch:b1:lewis-resonance-vsepr}

The ozone molecule, \ce{O3}, is three oxygen atoms in a bent chain. Draw
it with the rules of Book 1 and one bond comes out double, the other
single: one bond should be short, the other long. Yet microwave
spectroscopy finds the two \ce{O-O} bonds of ozone exactly equal, at
\qty{127.8}{pm}, between the double bond of \ce{O2} (\qty{120.8}{pm}) and
the single bond of hydrogen peroxide (\qty{147.5}{pm}). No single drawing
is right; two drawings together are. This chapter makes the Lewis model
precise --- \omterm{def:b1:lewis-resonance-vsepr:formal-charge}{formal charges}, octets that are incomplete or exceeded,
resonance between several structures --- and then uses it to predict
the shapes of molecules and the \omterm{def:b1:lewis-resonance-vsepr:dipole}{dipole moments} that follow from them.
% ledger: geo:O3.rOO, re:O2, geo:H2O2.rOO

\begin{recall}
Book 1 (grade 10) described the \omterm{def:b1:lewis-resonance-vsepr:lewis}{covalent bond} as a shared pair of
electrons, drew \omterm{def:b1:lewis-resonance-vsepr:lewis}{Lewis structures} with bonding and \omterm{def:b1:lewis-resonance-vsepr:lewis}{lone pairs}, and
predicted four shapes (linear, bent, trigonal, tetrahedral).
\omterm{def:b1:periodicity:electronegativity}{Electronegativity} and its differences were made quantitative in
\cref{ch:b1:periodicity}; the \omterm{def:b1:quantum-numbers:configuration}{valence electrons} of an atom are read from
its configuration (\cref{def:b1:quantum-numbers:configuration}).
\end{recall}

\section{Lewis structures and formal charges}

\begin{definition}[Covalent bond, Lewis structure]\label{def:b1:lewis-resonance-vsepr:lewis}
A \emph{covalent bond}\index{covalent bond} is a pair of electrons shared
by two atoms, the \emph{bonding pair}\index{bonding pair}; two or three
pairs make a double or a triple bond. A pair of \omterm{def:b1:quantum-numbers:configuration}{valence electrons} that
belongs to a single atom is a \emph{lone pair}\index{lone pair}. A
\emph{Lewis structure}\index{Lewis structure} shows every valence
electron of a molecule or ion, as bonds (lines between atoms) and lone
pairs (lines or pairs of dots on an atom). By the \emph{octet
rule}\index{octet rule}, the atoms of \omterm{def:b1:periodicity:table}{period} 2 (C, N, O, F) are
surrounded in a stable structure by four pairs, eight electrons; by the
\emph{duet rule}\index{duet rule}, hydrogen by one pair.
\end{definition}

\begin{definition}[Formal charge]\label{def:b1:lewis-resonance-vsepr:formal-charge}
The \emph{formal charge}\index{formal charge} of an atom in a \omterm{def:b1:lewis-resonance-vsepr:lewis}{Lewis
structure} is
\[
  q_F = N_v - N_{\text{lone}} - \tfrac12 N_{\text{bond}},
\]
where $N_v$ is the number of \omterm{def:b1:quantum-numbers:configuration}{valence electrons} of the free atom,
$N_{\text{lone}}$ the number of electrons in its \omterm{def:b1:lewis-resonance-vsepr:lewis}{lone pairs} and
$N_{\text{bond}}$ the number of electrons in the bonds it forms. It is
written $\oplus$ or $\ominus$ next to the atom when it is not zero.
\end{definition}

\begin{proposition}[Formal charges add up to the charge]\label{prop:b1:lewis-resonance-vsepr:formal-sum}
The sum of the \omterm{def:b1:lewis-resonance-vsepr:formal-charge}{formal charges} of the atoms of a \omterm{def:b1:lewis-resonance-vsepr:lewis}{Lewis structure} equals
the total charge of the molecule or ion.
\end{proposition}

\begin{proof}
Summing $q_F$ over the atoms gives $\sum N_v - (\text{electrons in lone
pairs} + \text{electrons in bonds})$, because each bond's electrons are
counted half on each of its two atoms. The bracket is the total number
of \omterm{def:b1:quantum-numbers:configuration}{valence electrons} drawn in the structure, which equals $\sum N_v$ minus
the charge $z$ of the species. Hence $\sum q_F = z$.
\end{proof}

\begin{method}[Drawing a Lewis structure]\label{met:b1:lewis-resonance-vsepr:lewis}
To draw the \omterm{def:b1:lewis-resonance-vsepr:lewis}{Lewis structure} of a molecule or ion:
\begin{enumerate}
  \item count the \omterm{def:b1:quantum-numbers:configuration}{valence electrons} $N = \sum N_v - z$ and the pairs $N/2$;
  \item draw the skeleton: the least electronegative atom (not H) is
        usually central; hydrogen and fluorine are always terminal;
  \item complete the octets of the terminal atoms with \omterm{def:b1:lewis-resonance-vsepr:lewis}{lone pairs}, then
        place the remaining pairs on the central atom;
  \item if the central atom lacks an octet, turn \omterm{def:b1:lewis-resonance-vsepr:lewis}{lone pairs} of its
        neighbours into multiple bonds;
  \item compute the \omterm{def:b1:lewis-resonance-vsepr:formal-charge}{formal charges}; among possible structures prefer the
        one with the fewest \omterm{def:b1:lewis-resonance-vsepr:formal-charge}{formal charges}, and with negative charges on
        the most electronegative atoms.
\end{enumerate}
\end{method}

\begin{example}[Nitric acid]\label{ex:b1:lewis-resonance-vsepr:hno3}
\ce{HNO3}: $N = 1 + 5 + 3 \times 6 = 24$ electrons, 12 pairs. Skeleton:
nitrogen central, bonded to three oxygens, one of which carries the
hydrogen. Completing octets with single bonds leaves nitrogen with six
electrons; one oxygen \omterm{def:b1:lewis-resonance-vsepr:lewis}{lone pair} becomes an \ce{N=O} bond. \omterm{def:b1:lewis-resonance-vsepr:formal-charge}{Formal charges}:
nitrogen $5 - 0 - 4 = +1$; the singly bonded oxygen without hydrogen
$6 - 6 - 1 = -1$; the others 0. The total is 0, as
\cref{prop:b1:lewis-resonance-vsepr:formal-sum} requires.
\end{example}

\begin{omfigure}
\setchemfig{atom sep=2.4em}
\begin{tikzpicture}[font=\small]
  \node at (0,0) {\large\chemfig{H-\charge{90=\|,270=\|}{O}-\charge{90:2pt=$\scriptstyle\oplus$}{N}(=[:45]\charge{0=\|,90=\|}{O})-[:-45]\charge{0=\|,240=\|,300=\|,45:3pt=$\scriptstyle\ominus$}{O}}};
  \node[below] at (0,-1.5) {nitric acid, \ce{HNO3}};
  \node at (5.2,0) {\large\chemfig{\charge{180=\|}{N}~\charge{0=\|}{N}}};
  \node[below] at (5.2,-1.5) {dinitrogen, \ce{N2}};
  \node at (9.6,0) {\large\chemfig{\charge{135=\|,225=\|,90:2pt=$\scriptstyle\ominus$}{N}=\charge{90:2pt=$\scriptstyle\oplus$}{N}=\charge{45=\|,315=\|}{O}}};
  \node[below] at (9.6,-1.5) {dinitrogen monoxide, \ce{N2O}};
\end{tikzpicture}

{\small \omterm{def:b1:lewis-resonance-vsepr:lewis}{Lewis structures} with \omterm{def:b1:lewis-resonance-vsepr:lewis}{lone pairs} and \omterm{def:b1:lewis-resonance-vsepr:formal-charge}{formal charges}.}
\end{omfigure}

\section{Beyond the octet}

\begin{definition}[Electron-deficient and hypervalent molecules]\label{def:b1:lewis-resonance-vsepr:hypervalent}
A molecule in which an atom has fewer than eight \omterm{def:b1:quantum-numbers:configuration}{valence electrons} is
\emph{electron-deficient}\index{electron-deficient}: in \ce{BF3} boron
has six. A molecule in which an atom of \omterm{def:b1:periodicity:table}{period} 3 or beyond is surrounded
by more than four pairs is \emph{hypervalent}\index{hypervalent}: in
\ce{PCl5} phosphorus forms five bonds, in \ce{SF6} sulfur six. Molecules
with an odd number of electrons, such as \ce{NO} and \ce{NO2}, cannot
satisfy the \omterm{def:b1:lewis-resonance-vsepr:lewis}{octet rule} on every atom: one electron is unpaired.
\end{definition}

\begin{remark}[Why period 3 and not period 2]\label{rem:b1:lewis-resonance-vsepr:hypervalent}
Nitrogen forms \ce{NF3} but never \ce{NF5}, while phosphorus forms both
\ce{PF3} and \ce{PF5}. A period-2 atom is too small to bond to more than
four neighbours; a period-3 atom is large enough to hold five or six.
Drawings of \ce{SO4^{2-}} or \ce{H2SO4} with sulfur double bonds (12
electrons around S) or with single bonds and \omterm{def:b1:lewis-resonance-vsepr:formal-charge}{formal charges} (an octet)
are both found; the second reflects better the charge distribution, and
both give the right geometry.
\end{remark}

\begin{definition}[Lewis acid, Lewis base, dative bond]\label{def:b1:lewis-resonance-vsepr:lewis-acid}
A \emph{Lewis acid}\index{Lewis acid} is a species able to accept a
pair of electrons into a vacant place of its valence shell (\ce{BF3},
\ce{AlCl3}, \ce{H+}, metal cations); a \emph{Lewis base}\index{Lewis base}
is a species with a \omterm{def:b1:lewis-resonance-vsepr:lewis}{lone pair} it can share (\ce{NH3}, \ce{H2O},
\ce{F-}). When the base gives the pair to the acid, the bond formed is a
\emph{dative bond}\index{dative bond}: a \omterm{def:b1:lewis-resonance-vsepr:lewis}{covalent bond} whose two
electrons both came from one partner.
\end{definition}

\begin{example}[Ammonia and boron trifluoride]\label{ex:b1:lewis-resonance-vsepr:adduct}
\ce{NH3 + BF3 -> H3NBF3}. In the adduct, boron completes its octet;
nitrogen, which now shares its \omterm{def:b1:lewis-resonance-vsepr:lewis}{lone pair}, carries the \omterm{def:b1:lewis-resonance-vsepr:formal-charge}{formal charge}
$+1$ and boron $-1$. Once formed, the \ce{N-B} bond is an ordinary
\omterm{def:b1:lewis-resonance-vsepr:lewis}{covalent bond}.
\end{example}

\begin{omfigure}
\setchemfig{atom sep=2.1em}
\begin{tikzpicture}[font=\small]
  \node at (0,0) {\large\chemfig{H-\charge{90=\|}{N}(-[:240]H)(-[:300]H)}};
  \node at (1.6,0) {$+$};
  \node at (3.1,0) {\large\chemfig{B(-[:90]F)(-[:210]F)(-[:330]F)}};
  \draw[-{Stealth}, thick] (4.3,0) -- (5.3,0);
  \node at (7.4,0) {\large\chemfig{H-\charge{45:1pt=$\scriptstyle\oplus$}{N}(-[:90]H)(-[:270]H)-\charge{45:1pt=$\scriptstyle\ominus$}{B}(-[:90]F)(-[:270]F)-F}};
\end{tikzpicture}

{\small A \omterm{def:b1:lewis-resonance-vsepr:lewis-acid}{dative bond}: the \omterm{def:b1:lewis-resonance-vsepr:lewis}{lone pair} of ammonia (a \omterm{def:b1:lewis-resonance-vsepr:lewis-acid}{Lewis base}) fills the
vacant place of boron trifluoride (a \omterm{def:b1:lewis-resonance-vsepr:lewis-acid}{Lewis acid}). \omterm{def:b1:lewis-resonance-vsepr:lewis}{Lone pairs} of fluorine
are not drawn.}
\end{omfigure}

\section{Resonance}

\begin{definition}[Resonance structures, resonance hybrid]\label{def:b1:lewis-resonance-vsepr:resonance}
When the electrons of a molecule can be placed in several \omterm{def:b1:lewis-resonance-vsepr:lewis}{Lewis
structures} that differ only by the position of $\pi$ bonds and \omterm{def:b1:lewis-resonance-vsepr:lewis}{lone
pairs} --- the nuclei staying in place --- each is a \emph{resonance
structure}\index{resonance structure} (or mesomeric form), linked to the
others by a double-headed arrow $\leftrightarrow$. The real molecule is
none of them: it is the \emph{resonance hybrid}\index{resonance hybrid},
a single structure whose electron distribution is a weighted average of
theirs.
\end{definition}

\begin{remark}[An arrow, not an equilibrium]\label{rem:b1:lewis-resonance-vsepr:arrow}
The arrow $\leftrightarrow$ does not describe a reaction: ozone does not
switch from one structure to the other. The hybrid is one molecule,
described by several drawings because one drawing is not enough. Never
use $\rightleftharpoons$ between \omterm{def:b1:lewis-resonance-vsepr:resonance}{resonance structures}.
\end{remark}

\begin{method}[Writing and weighting resonance structures]\label{met:b1:lewis-resonance-vsepr:resonance}
To find the \omterm{def:b1:lewis-resonance-vsepr:resonance}{resonance structures} of a species:
\begin{enumerate}
  \item start from one \omterm{def:b1:lewis-resonance-vsepr:lewis}{Lewis structure}; look for a \omterm{def:b1:lewis-resonance-vsepr:lewis}{lone pair} or a
        $\pi$ bond next to a $\pi$ bond, a vacant place or a positive
        charge;
  \item move pairs with curly arrows (a \omterm{def:b1:lewis-resonance-vsepr:lewis}{lone pair} becomes a $\pi$ bond, a
        $\pi$ bond becomes a \omterm{def:b1:lewis-resonance-vsepr:lewis}{lone pair}), never moving a nucleus and never
        exceeding the octet of a period-2 atom;
  \item recompute the \omterm{def:b1:lewis-resonance-vsepr:formal-charge}{formal charges};
  \item weight the structures: the most important ones obey the \omterm{def:b1:lewis-resonance-vsepr:lewis}{octet
        rule}, have the fewest \omterm{def:b1:lewis-resonance-vsepr:formal-charge}{formal charges}, and put negative charges on
        the most electronegative atoms; equivalent structures weigh the
        same.
\end{enumerate}
\end{method}

\begin{example}[Ozone, nitrate, benzene]\label{ex:b1:lewis-resonance-vsepr:examples}
Ozone has two equivalent structures: each \ce{O-O} bond is double in one
and single in the other, so both have a bond order $1.5$ and the same
length, \qty{127.8}{pm}, between \ce{O=O} and \ce{O-O}. The nitrate ion
\ce{NO3-} has three equivalent structures; each \ce{N-O} bond has order
$4/3$. Benzene, \ce{C6H6}, has two Kekulé structures; its six \ce{C-C}
bonds are equal, \qty{139.7}{pm}, between the single bond of ethane
(\qty{153.6}{pm}) and the double bond of ethene (\qty{133.9}{pm}).
% ledger: geo:O3.rOO, geo:C6H6.rCC, geo:C2H6.rCC, geo:C2H4.rCC
\end{example}

\begin{omfigure}
\vspace*{10pt}
\large
\setchemfig{atom sep=2.4em}
\schemestart
\chemfig{\charge{135=\|,225=\|}{O}=[:-30]\charge{270=\|,90:2pt=$\scriptstyle\oplus$}{O}-[:30]\charge{90=\|,0=\|,270=\|,45:5pt=$\scriptstyle\ominus$}{O}}
\arrow{<->}
\chemfig{\charge{90=\|,180=\|,270=\|,135:5pt=$\scriptstyle\ominus$}{O}-[:-30]\charge{270=\|,90:2pt=$\scriptstyle\oplus$}{O}=[:30]\charge{45=\|,315=\|}{O}}
\schemestop

\medskip
\schemestart
\chemfig{*6(=-=-=-)}
\arrow{<->}
\chemfig{*6(-=-=-=)}
\schemestop
\hspace{3em}
\chemfig{\charge{270:2pt=$\scriptstyle\oplus$}{N}(=[:90]\charge{45=\|,135=\|}{O})(-[:210]\charge{150=\|,240=\|,300=\|,90:2pt=$\scriptstyle\ominus$}{O})-[:330]\charge{30=\|,300=\|,240=\|,90:2pt=$\scriptstyle\ominus$}{O}}
\normalsize
\par\vspace{14pt}
{\small Resonance in ozone (one \omterm{def:b1:lewis-resonance-vsepr:lewis}{lone pair} of the right-hand oxygen
becomes a $\pi$ bond while the $\pi$ bond on the left becomes a \omterm{def:b1:lewis-resonance-vsepr:lewis}{lone
pair}) and in benzene, and one of
the three equivalent structures of the nitrate ion, whose charge is
shared by the three oxygens.}
\end{omfigure}

\begin{definition}[$\sigma$ bond, $\pi$ bond]\label{def:b1:lewis-resonance-vsepr:sigma-pi}
A single bond is a \emph{$\sigma$ bond}\index{sigma bond@$\sigma$ bond}:
its electron pair lies along the axis joining the two nuclei, from the
head-on overlap of two orbitals. The second and third bonds of a multiple
bond are \emph{$\pi$ bonds}\index{pi bond@$\pi$ bond}: their pairs lie on
either side of the axis, from the side-on overlap of two p orbitals
perpendicular to it. A double bond is one $\sigma$ and one $\pi$ bond; a
triple bond one $\sigma$ and two $\pi$.
\end{definition}

\begin{omfigure}
\begin{tikzpicture}[font=\small]
  % sigma: two lobes overlapping head-on
  \fill[omDef!35, opacity=0.8] (0,0) ellipse (0.9 and 0.42);
  \fill[omProp!35, opacity=0.8] (1.1,0) ellipse (0.9 and 0.42);
  \fill (0.1,0) circle (1.6pt); \fill (1.0,0) circle (1.6pt);
  \draw[dashed, black!50] (-1.2,0) -- (2.3,0);
  \node[below] at (0.55,-1.05) {$\sigma$: head-on, on the axis};
  % pi: two p orbitals side by side
  \begin{scope}[xshift=5.2cm]
    \foreach \x/\c in {0/omDef, 1.2/omProp} {
      \fill[\c!35, opacity=0.85] (\x,0.45) ellipse (0.32 and 0.45);
      \fill[\c!15, opacity=0.85] (\x,-0.45) ellipse (0.32 and 0.45);
      \draw[\c!70!black] (\x,0.45) ellipse (0.32 and 0.45) (\x,-0.45) ellipse (0.32 and 0.45);
      \fill (\x,0) circle (1.6pt);
    }
    \draw[dashed, black!50] (-0.7,0) -- (1.9,0);
    \draw[omThm, thick, densely dotted] (0.25,0.75) -- (0.95,0.75) (0.25,-0.75) -- (0.95,-0.75);
    \node[below] at (0.6,-1.05) {$\pi$: side-on, above and below};
  \end{scope}
\end{tikzpicture}

{\small The two kinds of overlap. A $\sigma$ bond allows rotation about
its axis; a $\pi$ bond does not, which is why a double bond is rigid
(\cref{ch:b1:stereochemistry-in-depth}).}
\end{omfigure}

\section{The VSEPR model}

\begin{definition}[VSEPR model, steric number]\label{def:b1:lewis-resonance-vsepr:vsepr}
In the \emph{VSEPR model}\index{VSEPR model} (valence-shell
electron-pair repulsion), the pairs around a central atom A repel one
another and take the arrangement that keeps them farthest apart. A
molecule is written \ce{AX_nE_m}, where $n$ is the number of atoms bonded
to A (a multiple bond counts once) and $m$ the number of \omterm{def:b1:lewis-resonance-vsepr:lewis}{lone pairs} of
A. The \emph{steric number}\index{steric number} $n + m$ fixes the
arrangement of the pairs: 2 linear, 3 trigonal planar, 4 tetrahedral, 5
trigonal bipyramidal, 6 octahedral. The shape of the molecule is the
arrangement of its atoms only.
\end{definition}

\begin{proposition}[VSEPR geometries]\label{prop:b1:lewis-resonance-vsepr:geometries}
The \omterm{def:b1:lewis-resonance-vsepr:vsepr}{VSEPR model} predicts the shapes and ideal angles of the table below;
measured angles are close.
\begin{center}
\small
\begin{tabular}{llllc}
\hline
type & arrangement of pairs & shape & example & measured angle \\
\hline
\ce{AX2} & linear & linear & \ce{CO2} & $180^\circ$ \\
\ce{AX3} & trigonal planar & trigonal planar & \ce{BF3} & $120^\circ$ \\
\ce{AX2E} & trigonal planar & bent & \ce{SO2} & $119.5^\circ$ \\
\ce{AX4} & tetrahedral & tetrahedral & \ce{CH4} & $109.5^\circ$ \\
\ce{AX3E} & tetrahedral & trigonal pyramidal & \ce{NH3} & $106.7^\circ$ \\
\ce{AX2E2} & tetrahedral & bent & \ce{H2O} & $104.5^\circ$ \\
\ce{AX5} & trigonal bipyramidal & trigonal bipyramidal & \ce{PCl5} & $90^\circ$, $120^\circ$ \\
\ce{AX4E} & trigonal bipyramidal & seesaw & \ce{SF4} & $101.6^\circ$, $173.1^\circ$ \\
\ce{AX3E2} & trigonal bipyramidal & T-shaped & \ce{ClF3} & $87.5^\circ$ \\
\ce{AX2E3} & trigonal bipyramidal & linear & \ce{XeF2} & $180^\circ$ \\
\ce{AX6} & octahedral & octahedral & \ce{SF6} & $90^\circ$ \\
\ce{AX5E} & octahedral & square pyramidal & \ce{BrF5} & --- \\
\ce{AX4E2} & octahedral & square planar & \ce{XeF4} & --- \\
\hline
\end{tabular}
\end{center}
% ledger: geo:CO2.aOCO, geo:BF3.aFBF, geo:SO2.aOSO, geo:CH4.aHCH,
% geo:NH3.aHNH, geo:H2O.aHOH, geo:SF6.aFSF, geo:SF4.aFSF2, geo:SF4.aFSF,
% geo:ClF3.aFClF, geo:XeF2.aFXeF
\end{proposition}
\admitted

\begin{proposition}[Lone pairs close the angles]\label{prop:b1:lewis-resonance-vsepr:lone-pair-angles}
A \omterm{def:b1:lewis-resonance-vsepr:lewis}{lone pair}, held by one nucleus only, spreads more than a \omterm{def:b1:lewis-resonance-vsepr:lewis}{bonding pair}
and repels more strongly: repulsions decrease in the order \omterm{def:b1:lewis-resonance-vsepr:lewis}{lone
pair}/\omterm{def:b1:lewis-resonance-vsepr:lewis}{lone pair} $>$ \omterm{def:b1:lewis-resonance-vsepr:lewis}{lone pair}/\omterm{def:b1:lewis-resonance-vsepr:lewis}{bonding pair} $>$ \omterm{def:b1:lewis-resonance-vsepr:lewis}{bonding pair}/\omterm{def:b1:lewis-resonance-vsepr:lewis}{bonding pair}.
Hence the angles between bonds shrink as \omterm{def:b1:lewis-resonance-vsepr:lewis}{lone pairs} are added
($109.5^\circ$ in \ce{CH4}, $106.7^\circ$ in \ce{NH3}, $104.5^\circ$ in
\ce{H2O}), and in a trigonal bipyramid \omterm{def:b1:lewis-resonance-vsepr:lewis}{lone pairs} take the equatorial
positions, where they have only two neighbours at $90^\circ$.
% ledger: geo:CH4.aHCH, geo:NH3.aHNH, geo:H2O.aHOH
\end{proposition}
\admitted

\begin{omfigure}
\setchemfig{atom sep=2.6em}
\renewcommand{\arraystretch}{1.2}
\begin{tabular}{ccccc}
\chemfig{O=C=O} &
\chemfig{B(-[:90]F)(-[:210]F)(-[:330]F)} &
\chemfig{C(-[:90]H)(-[:200]H)(<[:280]H)(<:[:335]H)} &
\chemfig{P(-[:90]Cl)(-[:270]Cl)(-[:180]Cl)(<[:320]Cl)(<:[:30]Cl)} &
\chemfig{S(-[:90]F)(-[:270]F)(-[:180]F)(-[:0]F)(<[:225]F)(<:[:45]F)} \\
\small \ce{AX2}, \ce{CO2} & \small \ce{AX3}, \ce{BF3} & \small \ce{AX4}, \ce{CH4} &
\small \ce{AX5}, \ce{PCl5} & \small \ce{AX6}, \ce{SF6} \\[16pt]
\chemfig{\charge{90=\:}{S}(=[:210]O)(=[:330]O)} &
\chemfig{\charge{90=\:}{N}(-[:200]H)(<[:280]H)(<:[:335]H)} &
\chemfig{\charge{60=\:,120=\:}{O}(-[:215]H)(-[:325]H)} &
\chemfig{\charge{0=\:}{S}(-[:90]F)(-[:270]F)(<[:210]F)(<:[:150]F)} &
\chemfig{Xe(-[:0]F)(-[:90]F)(-[:180]F)(-[:270]F)} \\
\small \ce{AX2E}, \ce{SO2} & \small \ce{AX3E}, \ce{NH3} & \small \ce{AX2E2}, \ce{H2O} &
\small \ce{AX4E}, \ce{SF4} & \small \ce{AX4E2}, \ce{XeF4} \\
\end{tabular}

\medskip
{\small Shapes predicted by VSEPR. A wedge points towards the reader, a
hashed wedge away; plain bonds lie in the plane of the page. \omterm{def:b1:lewis-resonance-vsepr:lewis}{Lone pairs}
of the central atom are shown as pairs of dots; \ce{XeF4} is drawn face
on, its two \omterm{def:b1:lewis-resonance-vsepr:lewis}{lone pairs} pointing towards and away from the reader.}
\end{omfigure}

\begin{method}[Predicting a shape]\label{met:b1:lewis-resonance-vsepr:vsepr}
To predict the shape around a central atom:
\begin{enumerate}
  \item draw the \omterm{def:b1:lewis-resonance-vsepr:lewis}{Lewis structure} (\cref{met:b1:lewis-resonance-vsepr:lewis});
  \item count the atoms bonded to the central atom ($n$) and its \omterm{def:b1:lewis-resonance-vsepr:lewis}{lone
        pairs} ($m$); write \ce{AX_nE_m};
  \item read the arrangement from $n + m$ and the shape from the
        positions of the $n$ atoms;
  \item correct the ideal angles: \omterm{def:b1:lewis-resonance-vsepr:lewis}{lone pairs} and multiple bonds take more
        room than single bonds.
\end{enumerate}
\end{method}

\begin{definition}[Hybridisation]\label{def:b1:lewis-resonance-vsepr:hybridisation}
The \emph{hybridisation}\index{hybridisation} of an atom is a
description of its bonds by orbitals pointing along the directions
predicted by VSEPR: an atom of \omterm{def:b1:lewis-resonance-vsepr:vsepr}{steric number} 4 is said to be sp$^3$
(four equivalent orbitals at $109.5^\circ$), of \omterm{def:b1:lewis-resonance-vsepr:vsepr}{steric number} 3 sp$^2$
(three orbitals at $120^\circ$ in a plane, one p orbital left
perpendicular to it), of \omterm{def:b1:lewis-resonance-vsepr:vsepr}{steric number} 2 sp (two orbitals at $180^\circ$,
two p orbitals left). The p orbitals left over form the $\pi$ bonds.
\end{definition}

\begin{example}[The carbons of organic chemistry]\label{ex:b1:lewis-resonance-vsepr:carbon}
In ethane each carbon is sp$^3$, tetrahedral; in ethene sp$^2$, the six
atoms in one plane with angles near $120^\circ$ (measured \ce{H-C-H}
$117.6^\circ$); in ethyne sp, the four atoms on one line. The bond
lengths shrink from 153.6 to 133.9 to \qty{120.3}{pm} as $\pi$ bonds are
added.
% ledger: geo:C2H4.aHCH, geo:C2H6.rCC, geo:C2H4.rCC, geo:C2H2.rCC
\end{example}

\begin{remark}[A description, not a cause]\label{rem:b1:lewis-resonance-vsepr:hybrid-limit}
\omterm{def:b1:lewis-resonance-vsepr:hybridisation}{Hybridisation} describes a geometry already known; it does not explain
it, and it fails for the energies of electrons, which molecular orbital
theory, in the Year 2 volume, accounts for. It remains the everyday
language of organic chemistry: an ``sp$^2$ carbon'' means a planar
carbon with a p orbital free for a $\pi$ bond.
\end{remark}

\section{Polar molecules: dipole moments}

\begin{definition}[Bond dipole, dipole moment]\label{def:b1:lewis-resonance-vsepr:dipole}
In a bond between atoms of different electronegativities the electrons
lean towards the more electronegative atom, which bears a partial charge
$-\delta e$ while its partner bears $+\delta e$ ($0 < \delta < 1$). The
\emph{bond dipole}\index{bond dipole} is the vector $\vec\mu = \delta e\,
\vec d$, of norm $\delta e d$, pointing from the negative to the positive
partial charge, $\vec d$ being the vector between them. The
\emph{dipole moment}\index{dipole moment} of a molecule is the vector sum
of its bond dipoles (and of the contributions of its \omterm{def:b1:lewis-resonance-vsepr:lewis}{lone pairs}); it is
expressed in coulomb metres or in debyes, $\qty{1}{\debye} =
\qty{3.336e-30}{C.m}$. A molecule with a non-zero dipole moment is a
\emph{polar molecule}\index{polar molecule}.
% ledger: const:c (1 D = 1e-21/c C.m)
\end{definition}

\begin{remark}[Direction convention]\label{rem:b1:lewis-resonance-vsepr:convention}
Physics draws the dipole vector from the negative to the positive
charge, the convention used here. Older chemistry books draw an arrow
with a crossed tail pointing towards the negative end; only the
direction differs.
\end{remark}

\begin{proposition}[Dipole of a symmetric molecule]\label{prop:b1:lewis-resonance-vsepr:dipole-sum}
A molecule \ce{AX_n} without \omterm{def:b1:lewis-resonance-vsepr:lewis}{lone pairs} on A (linear \ce{AX2}, trigonal
\ce{AX3}, tetrahedral \ce{AX4}, \dots) has a zero \omterm{def:b1:lewis-resonance-vsepr:dipole}{dipole moment}, whatever
the polarity of its bonds. A bent \ce{AX2} molecule of angle $\theta$
whose bonds have the dipole $\mu_b$ has the \omterm{def:b1:lewis-resonance-vsepr:dipole}{dipole moment}
\[
  \mu = 2\mu_b \cos\frac{\theta}{2} .
\]
\end{proposition}

\begin{proof}
In the symmetric shapes the bond vectors are equal in norm and their
directions are arranged symmetrically around A, so they add up to zero
(for \ce{AX2} linear they are opposite; for \ce{AX3} at $120^\circ$
three vectors of the same norm sum to zero; for \ce{AX4} the sum is
invariant under the rotations of the tetrahedron, so it is zero). For the
bent molecule, each bond vector makes the angle $\theta/2$ with the
bisector; the components perpendicular to the bisector cancel and those
along it add: $2\mu_b\cos(\theta/2)$.
\end{proof}

\begin{omfigure}
\begin{tikzpicture}[font=\small, >={Stealth}]
  % water
  \node[aO] (o) at (0,0) {};
  \node[aH] (h1) at (-0.95,-0.73) {};
  \node[aH] (h2) at (0.95,-0.73) {};
  \begin{scope}[on background layer]
    \draw[bond] (o.center) -- (h1.center); \draw[bond] (o.center) -- (h2.center);
  \end{scope}
  \draw[->, omProp, thick] (-0.25,0.15) -- (-0.95,-0.38);
  \draw[->, omProp, thick] (0.25,0.15) -- (0.95,-0.38);
  \draw[->, omThm, very thick] (0,-0.25) -- (0,-1.45);
  \node[omThm, right] at (0.05,-1.25) {$\vec\mu$};
  \node at (0,0.55) {$\delta^-$};
  \node at (-1.25,-1.05) {$\delta^+$}; \node at (1.25,-1.05) {$\delta^+$};
  \node[below] at (0,-1.7) {\ce{H2O}: $\mu = \qty{1.86}{\debye}$};
  % carbon dioxide
  \begin{scope}[xshift=5cm]
    \node[aO] (a) at (-1.1,0) {}; \node[aC] (c) at (0,0) {}; \node[aO] (b) at (1.1,0) {};
    \begin{scope}[on background layer]
      \foreach \d in {-1.6,1.6} {
        \draw[bond, line width=1.1pt] ([yshift=\d pt]a.center) -- ([yshift=\d pt]c.center);
        \draw[bond, line width=1.1pt] ([yshift=\d pt]c.center) -- ([yshift=\d pt]b.center);}
    \end{scope}
    \draw[->, omProp, thick] (-0.75,0.45) -- (-0.2,0.45);
    \draw[->, omProp, thick] (0.75,0.45) -- (0.2,0.45);
    \node[below] at (0,-1.7) {\ce{CO2}: bond dipoles cancel, $\mu = 0$};
  \end{scope}
\end{tikzpicture}

{\small \omterm{def:b1:lewis-resonance-vsepr:dipole}{Bond dipoles} (orange, from the partial negative to the partial
positive charge) and their sum (red). In bent water they add; in linear
carbon dioxide they cancel.}
% ledger: dip:H2O
\end{omfigure}

\begin{example}[Polar and non-polar molecules]\label{ex:b1:lewis-resonance-vsepr:dipoles}
\ce{CO2}, \ce{BF3}, \ce{CH4} and \ce{CCl4} have zero \omterm{def:b1:lewis-resonance-vsepr:dipole}{dipole moments}.
\ce{H2O} (\qty{1.86}{\debye}), \ce{NH3} (\qty{1.48}{\debye}) and
\ce{SO2} (\qty{1.63}{\debye}) are polar. Along \ce{CH3Cl}, \ce{CH2Cl2},
\ce{CHCl3} the moment falls, 1.87, 1.62, \qty{1.04}{\debye}, to vanish in
\ce{CCl4}: the \ce{C-Cl} \omterm{def:b1:lewis-resonance-vsepr:dipole}{bond dipoles} increasingly cancel. Among the
hydrogen halides it follows the \omterm{def:b1:periodicity:electronegativity}{electronegativity} difference: HF 1.83,
HCl 1.09, HBr 0.83, HI \qty{0.45}{\debye}.
% ledger: dip:H2O, dip:NH3, dip:SO2, dip:CH3Cl, dip:CH2Cl2, dip:CHCl3,
% dip:CCl4, dip:HF, dip:HCl, dip:HBr, dip:HI
\end{example}

\section{Exercises}

\begin{exercise}[$\star$]\label{exo:b1:lewis-resonance-vsepr:1}
Draw the \omterm{def:b1:lewis-resonance-vsepr:lewis}{Lewis structures} of \ce{H2O}, \ce{NH3}, \ce{CO2}, \ce{HCN},
\ce{CH2O} (methanal) and \ce{NH4+}, and give the \omterm{def:b1:lewis-resonance-vsepr:formal-charge}{formal charges}.
\end{exercise}

\begin{exercise}[$\star$]\label{exo:b1:lewis-resonance-vsepr:2}
Give the type \ce{AX_nE_m}, the shape and the approximate angle of
\ce{H2S}, \ce{PCl3}, \ce{BF3}, \ce{SiH4} and \ce{CO3^{2-}} (central atom
first in each formula).
\end{exercise}

\begin{exercise}[$\star$]\label{exo:b1:lewis-resonance-vsepr:3}
Which of these molecules are polar: \ce{CCl4}, \ce{CH2Cl2}, \ce{BF3},
\ce{NH3}, \ce{CO2}, \ce{SO2}? Justify from the shapes.
\end{exercise}

\begin{exercise}[$\star$]\label{exo:b1:lewis-resonance-vsepr:4}
Identify the \omterm{def:b1:lewis-resonance-vsepr:lewis-acid}{Lewis acid} and the \omterm{def:b1:lewis-resonance-vsepr:lewis-acid}{Lewis base} in \ce{H+ + H2O -> H3O+} and
in \ce{Cu^{2+} + 4NH3 -> [Cu(NH3)4]^{2+}}. Which bonds are dative?
\end{exercise}

\begin{exercise}[$\star\star$]\label{exo:b1:lewis-resonance-vsepr:5}
Write two \omterm{def:b1:lewis-resonance-vsepr:resonance}{resonance structures} of the ethanoate ion \ce{CH3COO-}. What
do they predict for the two \ce{C-O} bond lengths? Compare with methanoic
acid \ce{HCOOH}, whose two \ce{C-O} bonds measure 120.2 and
\qty{134.3}{pm}.
% ledger: geo:H2CO2.rCO, geo:H2CO2.rCO2
\end{exercise}

\begin{exercise}[$\star\star$]\label{exo:b1:lewis-resonance-vsepr:6}
Draw the \omterm{def:b1:lewis-resonance-vsepr:lewis}{Lewis structure} of the sulfate ion \ce{SO4^{2-}} with four
single \ce{S-O} bonds, compute the \omterm{def:b1:lewis-resonance-vsepr:formal-charge}{formal charges}, and predict the
shape. How many \omterm{def:b1:lewis-resonance-vsepr:resonance}{resonance structures} with two \ce{S=O} bonds and no
charge on sulfur can be written?
\end{exercise}

\begin{exercise}[$\star\star$]\label{exo:b1:lewis-resonance-vsepr:7}
Use VSEPR to predict the shapes of \ce{SF4}, \ce{ClF3}, \ce{XeF2} and
\ce{XeF4}. Explain why the \omterm{def:b1:lewis-resonance-vsepr:lewis}{lone pairs} of a trigonal bipyramid sit in the
equatorial plane.
\end{exercise}

\begin{exercise}[$\star\star$]\label{exo:b1:lewis-resonance-vsepr:8}
The angles \ce{H-X-H} are $104.5^\circ$ in \ce{H2O} and $92.1^\circ$ in
\ce{H2S}; $106.7^\circ$ in \ce{NH3} and $93.3^\circ$ in \ce{PH3}. Propose
an explanation in terms of the size and \omterm{def:b1:periodicity:electronegativity}{electronegativity} of the central
atom.
% ledger: geo:H2O.aHOH, geo:H2S.aHSH, geo:NH3.aHNH, geo:PH3.aHPH
\end{exercise}

\begin{exercise}[$\star\star$]\label{exo:b1:lewis-resonance-vsepr:9}
Give the \omterm{def:b1:lewis-resonance-vsepr:hybridisation}{hybridisation} of each carbon and of the nitrogen in
\ce{CH3-C#N} (ethanenitrile) and in \ce{CH2=CH-CH3}, and the number of
$\sigma$ and $\pi$ bonds in each molecule.
\end{exercise}

\begin{exercise}[$\star\star\star$]\label{exo:b1:lewis-resonance-vsepr:10}
The \omterm{def:b1:lewis-resonance-vsepr:dipole}{dipole moment} of \ce{NF3} is only \qty{0.24}{\debye}, against
\qty{1.48}{\debye} for \ce{NH3}, although the \ce{N-F} bond is more polar
than the \ce{N-H} bond. Using the \omterm{def:b1:lewis-resonance-vsepr:dipole}{bond dipoles} and the \omterm{def:b1:lewis-resonance-vsepr:lewis}{lone pair} of
nitrogen, explain the difference.
% ledger: dip:NF3, dip:NH3
\end{exercise}

\begin{exercise}[$\star\star\star$]\label{exo:b1:lewis-resonance-vsepr:11}
The \ce{SO2} molecule has an angle of $119.5^\circ$ and a \omterm{def:b1:lewis-resonance-vsepr:dipole}{dipole moment}
of \qty{1.63}{\debye}. Compute the dipole of one \ce{S-O} bond; with the
bond length \qty{143.2}{pm}, deduce the partial charge $\delta$ on each
oxygen.
% ledger: geo:SO2.aOSO, geo:SO2.rSO, dip:SO2
\end{exercise}

\begin{exercise}[$\star\star\star$]\label{exo:b1:lewis-resonance-vsepr:12}
Draw the most important \omterm{def:b1:lewis-resonance-vsepr:lewis}{Lewis structures} of dinitrogen monoxide \ce{N2O}
(linear, \ce{N-N-O}) and compute the \omterm{def:b1:lewis-resonance-vsepr:formal-charge}{formal charges}. The measured bond
lengths are \ce{N-N} \qty{112.8}{pm} and \ce{N-O} \qty{118.4}{pm}; in
\ce{N2} the triple bond measures \qty{109.8}{pm}. Which structure weighs
more?
% ledger: geo:N2O.rNN, geo:N2O.rNO, re:N2
\end{exercise}

\section{Problem: The Molecules of a Lightning Strike}

\begin{problem}[{Weekend problem --- the nitrogen and oxygen species made
when lightning heats air: Lewis structures, resonance, shapes, and the
partial charges of water}]
\label{pb:b1:lewis-resonance-vsepr:1}
A lightning bolt heats air to thousands of degrees: \ce{N2} and \ce{O2}
give \ce{NO}, then \ce{NO2}, \ce{N2O}, ozone and finally nitric acid in
rain. Data: $\qty{1}{\debye} = \qty{3.336e-30}{C.m}$, $e =
\qty{1.602e-19}{C}$. Measured values are given where needed.

\medskip
\noindent\textbf{Part I --- \omterm{def:b1:lewis-resonance-vsepr:lewis}{Lewis structures}.}
\begin{enumerate}
  \item Count the \omterm{def:b1:quantum-numbers:configuration}{valence electrons} of \ce{N2}, \ce{NO}, \ce{NO2},
        \ce{N2O}, \ce{O3} and \ce{HNO3}.
  \item Draw the \omterm{def:b1:lewis-resonance-vsepr:lewis}{Lewis structure} of \ce{N2}.
  \item Draw a \omterm{def:b1:lewis-resonance-vsepr:lewis}{Lewis structure} of \ce{NO}. Why can the \omterm{def:b1:lewis-resonance-vsepr:lewis}{octet rule} not be
        satisfied on both atoms?
  \item Draw two \omterm{def:b1:lewis-resonance-vsepr:lewis}{Lewis structures} of \ce{N2O}, with \ce{N=N=O} and with
        \ce{N#N-O}, and give the \omterm{def:b1:lewis-resonance-vsepr:formal-charge}{formal charges}.
  \item Draw a \omterm{def:b1:lewis-resonance-vsepr:lewis}{Lewis structure} of ozone and give the \omterm{def:b1:lewis-resonance-vsepr:formal-charge}{formal charges}.
  \item Draw a \omterm{def:b1:lewis-resonance-vsepr:lewis}{Lewis structure} of \ce{HNO3} and give the \omterm{def:b1:lewis-resonance-vsepr:formal-charge}{formal charges}.
  \item Which of the six species have an \omterm{def:b1:quantum-numbers:unpaired}{unpaired electron}?
\end{enumerate}

\medskip
\noindent\textbf{Part II --- Resonance and bond lengths.}
\begin{enumerate}[resume]
  \item Write the two \omterm{def:b1:lewis-resonance-vsepr:resonance}{resonance structures} of ozone. What bond order do
        they predict?
  \item The \ce{O-O} bond measures \qty{127.8}{pm} in ozone,
        \qty{120.8}{pm} in \ce{O2} and \qty{147.5}{pm} in \ce{H2O2}. Is
        the prediction confirmed?
  \item Write the three \omterm{def:b1:lewis-resonance-vsepr:resonance}{resonance structures} of \ce{NO3-} and give the
        order of each \ce{N-O} bond.
  \item In \ce{HNO3} the three \ce{N-O} bonds measure 140.6, 121.1 and
        \qty{119.9}{pm}. Assign them and explain using resonance.
  \item Explain why the angle \ce{O-N-O} of \ce{NO2} ($134^\circ$) is
        larger than $120^\circ$.
  \item Which of the two structures of question 4 places the negative
        \omterm{def:b1:lewis-resonance-vsepr:formal-charge}{formal charge} on the more electronegative atom?
\end{enumerate}

\medskip
\noindent\textbf{Part III --- Shapes and polarity.}
\begin{enumerate}[resume]
  \item Give the type \ce{AX_nE_m} of the central atom of \ce{O3},
        \ce{N2O}, \ce{NO3-} and the nitrogen of \ce{HNO3}, and their shapes.
  \item The ozone angle is $116.8^\circ$. Explain why it is below
        $120^\circ$.
  \item Ozone has a \omterm{def:b1:lewis-resonance-vsepr:dipole}{dipole moment} of \qty{0.53}{\debye}, \ce{N2O}
        \qty{0.16}{\debye}, \ce{N2} none. Explain.
  \item Why is \ce{CO2} non-polar while \ce{SO2} is polar?
  \item Predict the angle of \ce{NH3} relative to $109.5^\circ$, then
        compare with the measured $106.7^\circ$.
  \item Same question for water ($104.5^\circ$).
\end{enumerate}
% ledger: geo:O3.rOO, re:O2, geo:H2O2.rOO, geo:HNO3.rNO, geo:HNO3.rNO2,
% geo:HNO3.rNO3, geo:NO2.aONO, geo:O3.aOOO, dip:O3, dip:N2O

\medskip
\noindent\textbf{Part IV --- The partial charges of water.}
The \omterm{def:b1:lewis-resonance-vsepr:dipole}{dipole moment} of water is \qty{1.857}{\debye}, its angle
$104.48^\circ$ and its \ce{O-H} bond length \qty{95.8}{pm}.
% ledger: dip:H2O, geo:H2O.aHOH, geo:H2O.rOH
\begin{enumerate}[resume]
  \item Express the \omterm{def:b1:lewis-resonance-vsepr:dipole}{dipole moment} of water as a function of the \omterm{def:b1:lewis-resonance-vsepr:dipole}{bond
        dipole} $\mu_{OH}$ and the angle.
  \item Compute $\mu_{OH}$ in debyes.
  \item Convert it into coulomb metres.
  \item If the \omterm{def:b1:lewis-resonance-vsepr:dipole}{bond dipole} were made of two full charges $\pm e$ at the
        two nuclei, what would it be?
  \item Deduce the fractional ionic character $\delta$ of the \ce{O-H}
        bond, the ratio of the measured \omterm{def:b1:lewis-resonance-vsepr:dipole}{bond dipole} to that of full
        charges.
\end{enumerate}
\end{problem}
